\documentclass[12pt,a4paper]{article}
        \makeatletter
        \def\input@path{{../}}
        \makeatother
        \usepackage{almeria-fichas}

        \FichaMeta{Sesion 10}{Rutas y distancias por Almeria}{Bloque 2 · Geometria urbana}{Medir distancias en el plano y comparar rutas para elegir recorridos turisticos eficientes.}{Procedimientos}{Aplicar, Analizar, Crear}

        \begin{document}

        \FichaCabecera

        \begin{materialesbox}
        \begin{itemize}
    \item Ficha impresa
    \item Lapiz
    \item Regla
    \item Calculadora
\end{itemize}
        \end{materialesbox}

        \begin{pasosbox}
        \begin{enumerate}
    \item Localiza bien los puntos de partida y llegada.
    \item Si usas la formula de distancia, escribe todos los pasos.
    \item Convierte a metros o kilometros si la escala lo pide.
    \item Al final, interpreta el resultado con una frase.
\end{enumerate}
        \end{pasosbox}

        \begin{recordatoriobox}
        \begin{itemize}
    \item Distancia en el plano: $d=\sqrt{(x_2-x_1)^2+(y_2-y_1)^2}$.
    \item Si la escala es 1 unidad = 200 m, hay que convertir.
    \item La ruta mas corta no siempre coincide con la carretera real.
\end{itemize}
        \end{recordatoriobox}

        \begin{apoyobox}
        \begin{itemize}
    \item Haz primero las diferencias horizontal y vertical.
    \item Recuerda que sin raiz cuadrada tendrias $d^2$, no $d$.
    \item Escribe la escala al lado del calculo.
\end{itemize}
        \end{apoyobox}

        \BloomSection{aplicar}

\BloomExercise{aplicar}{1}{Distancias basicas}{Calcula la distancia entre $P(0,0)$ y $Q(3,4)$. Despues entre $R(1,2)$ y $S(4,6)$.}{6}

\BloomExercise{aplicar}{2}{Con escala}{Si en el plano 1 unidad son $200\,m$, convierte a metros la distancia entre $A(2,1)$ y $B(6,4)$.}{6}

\BloomSection{analizar}

\BloomExercise{analizar}{1}{Recta o carretera}{La distancia recta entre dos puntos es $8\,km$, pero la carretera mide $12\,km$. Analiza por que no coinciden.}{7}

\BloomExercise{analizar}{2}{Eleccion de trayecto}{Compara estas dos rutas: una con menos kilometros y otra con mas paradas utiles. Analiza que criterio te parece mejor para un turista.}{7}

\BloomSection{crear}

\BloomExercise{crear}{1}{Diseño propio}{Crea una produccion nueva relacionada con \emph{Rutas y distancias por Almeria}: un problema, una representacion, una mini guia, una explicacion visual o una propuesta de mejora.}{9}

\BloomExercise{crear}{2}{Alternativa mejorada}{Inventa una alternativa mas clara, mas util o mas creativa para trabajar este contenido. Debe estar alineada con el objetivo de la sesion y con la dimension procedimientos.}{9}

        \begin{evidenciabox}
        \begin{itemize}
    \item Calcular distancias entre puntos del plano.
    \item Interpretar la diferencia entre linea recta y ruta real.
    \item Justificar una eleccion de recorrido.
\end{itemize}
        \end{evidenciabox}

        \end{document}
